% Short version only: the voucher details that the body summarises. Text
% moved from the vouchers section to make room for the predicate boxes
% and algorithms; the full version carries the same material in its
% receive section and in Appendix~\ref{app:claim}.
\section{Vouchers in Detail}\label{app:vouchers}

\subsection{Cumulative Vouchers}\label{app:cumulative}

Every voucher against a source draws on the same leaf, and the first
redemption voids the rest. An issuer that pays one recipient repeatedly
therefore signs a \emph{cumulative} voucher for the running total against
the same claim leaf, and re-signs it for a larger amount at each payment:
2, then 5, then 9, one signature per payment. The recipient keeps only
the latest and, if the issuer stops cooperating, redeems that one alone;
a smaller one is never worth more to it, so no contest window.

Settled credit is retired by advancing $\mathsf{seq}$. Once the issuer
pays the recipient another way, it advances the source's $\mathsf{seq}$, which voids every outstanding
voucher against it: one \tr{UPDATE} per source leaf ($1{,}030\vB$), or
one \tr{REWRITE}$_h$ for a cohort of them (about $148\vB$ per leaf at
$h = 6$). Until then the issuer tracks every voucher it has signed as a
liability. Plain \tr{CLAIM} leaves a hole: anyone holding an old voucher
can redeem it, and the issuer holds them all, so redeeming the voucher
for 2 voids the one for 9. The credit leaf below closes it.

\subsection{The Dedicated Credit Leaf}\label{app:creditleaf}

% Shared by both versions: the dedicated credit leaf. Source: spec
% 7.5.3 and 10.3 (0.11.0). The heading and label are the includer's:
% a run-in paragraph of the voucher section (short) or of the receive
% section (full-receive.tex), so the text starts without a blank line.
A dedicated credit leaf is a source leaf that the issuer owns and whose
exit key is one recipient's, in a tree whose taptree carries
\tr{CLAIM}$^{x}$, \tr{UPDATE}$^{x}$ and the exact-payout exit, carries
\tr{REWRITE}$_h^{x}$ in place of \tr{REWRITE}$_h$ (or no rewrite at all),
and omits \tr{SPLIT}. \tr{CLAIM}$^{x}$ is \tr{CLAIM} with one further
signature over the transaction, by the source leaf's exit key: 7 bytes
of script and 64 of witness in FULL, the only tier that builds it. The
recipient's own redemptions carry that signature at no cost, and the
issuer cannot carve the leaf at all, not even with a voucher it signs to
itself. Over-issuance then has no second victim, since every voucher
pays the same recipient, who can watch the leaf's amount. Voiding needs
the leaf to change, and every transition that could change it needs the
exit key: \tr{CLAIM}$^{x}$, \tr{UPDATE}$^{x}$ and \tr{REWRITE}$_h^{x}$
by rule, while the exact exit and an eviction pay the whole balance to
that key and so void nothing the recipient is owed. The balance can only grow or pay its recipient, whose remaining
exposure is the issuer's liveness: a provider that stops issuing.

Three details keep it that way. The recipient checks from the taptree
root $m$ that no plain \tr{REWRITE}$_h$ is present, since it needs only
the owner's and the operator's signatures, both the issuer's. The issuer
funds the leaf only with credit it has issued or is about to, since any
excess is the recipient's on every path the issuer can take alone, and
tops it up by an amount-only increase under \tr{REWRITE}$_h^{x}$, which
needs only the operator's signature. And the leaf has a finite lease
with the operator as evictor, not a $\NUMS$ evictor, which every carve
would inherit, blocking the sweep of whatever cohort the claim leaf lands
in. The credit is one-way: if the recipient could also spend from it, an
older, larger voucher would be a stale state needing a window to
supersede, as in a bidirectional channel. To spend, the recipient
redeems into a channel leaf (\S\ref{sec:channels}).


\subsection{Hash-Locked Vouchers and Swaps}\label{app:hashlock}

\tr{CLAIM}$^{h}$ adds a payment hash to the signed message, under its
own tag so that the lock cannot be stripped, and requires the 32-byte
preimage: the script hashes it into the message, so a wrong preimage
fails as a bad signature. Paired with an HTLC or invoice on the same hash,
it is an ordered exchange, not an atomic swap: the issuer can void the
voucher until the redemption confirms, and script cannot give a voucher
a deadline. Five rules (spec~9.6.2) order it.
(1) The party the HTLC pays draws the preimage. (2) The voucher resolves
second; no HTLC waits on a preimage that only a redemption would reveal.
(3) The issuer leaves the source's $\mathsf{seq}$ alone from the moment
an HTLC on the hash is offered to it until the voucher is redeemed or
the HTLC fails. (4) No script enforces rule~3, so the holder redeems as
soon as it learns the preimage and bears the issuer risk until that
confirms. (5) The holder first checks that the source's lease outlasts
the HTLC's expiry plus $\lat(q)$. Usually the HTLC settles off-chain, so
such a swap costs at most one \tr{CLAIM}$^{h}$; it is also how a newcomer
buys a leaf over Lightning.
